% ============================================================
%  EthosNews — Technical & Product Documentation
%  Economic Times GenAI Hackathon 2026
%  Compile with: pdflatex (Overleaf default engine)
%
%  FILES NEEDED IN OVERLEAF (upload alongside this .tex):
%    fig_arch.png       — System architecture diagram
%    fig_ingest.png     — Ingestion pipeline diagram
%    fig_pnc.png        — PNC to feed flow diagram
%    img_factcheck.png  — Cropped infographic: fact-check section
%    img_clustering.png — Cropped infographic: clustering section
% ============================================================

\documentclass[9pt,a4paper,twocolumn]{article}

% ── Packages ────────────────────────────────────────────────
\usepackage[a4paper, top=16mm, bottom=16mm, left=13mm, right=13mm,
            columnsep=6mm]{geometry}
\usepackage{fontenc}
\usepackage{inputenc}
\usepackage[T1]{fontenc}
\usepackage{microtype}
\usepackage{graphicx}
\usepackage{xcolor}
\usepackage{booktabs}
\usepackage{tabularx}
\usepackage{array}
\usepackage{enumitem}
\usepackage{titlesec}
\usepackage{titling}
\usepackage{fancyhdr}
\usepackage{parskip}
\usepackage{amsmath}
\usepackage{mdframed}
\usepackage{multicol}
\usepackage{caption}
\usepackage{float}
\usepackage{lmodern}
\usepackage{hyperref}

% ── Broadsheet Colour Palette (from frontend globals.css) ───
\definecolor{paper}{HTML}{f7f2e9}
\definecolor{ink}{HTML}{0e0c09}
\definecolor{inksoft}{HTML}{2d2820}
\definecolor{inkmuted}{HTML}{7a6f63}
\definecolor{rule}{HTML}{c8bfb0}
\definecolor{etred}{HTML}{c8281e}
\definecolor{etgold}{HTML}{b5830a}
\definecolor{steel}{HTML}{2a4a8a}
\definecolor{mint}{HTML}{1a7a52}
\definecolor{reddim}{HTML}{f5e8e7}

% ── Page background ─────────────────────────────────────────
\pagecolor{paper}
\color{ink}

% ── Hyperlinks ───────────────────────────────────────────────
\hypersetup{
  colorlinks=true,
  linkcolor=steel,
  urlcolor=steel,
  citecolor=etgold,
}

% ── Section heading styles ──────────────────────────────────
\titleformat{\section}
  {\fontfamily{lmr}\selectfont\bfseries\large\color{ink}}
  {}
  {0pt}
  {\colorbox{ink}{\color{paper}\hspace{4pt}#1\hspace{4pt}}\vspace{2pt}\newline\textcolor{etred}{\rule{\columnwidth}{1.5pt}}}

\titleformat{\subsection}
  {\fontfamily{lmr}\selectfont\bfseries\normalsize\color{etred}}
  {}
  {0pt}
  {#1}

\titleformat{\subsubsection}
  {\fontfamily{lmr}\selectfont\bfseries\small\color{etgold}}
  {}
  {0pt}
  {#1}

\titlespacing{\section}{0pt}{8pt}{4pt}
\titlespacing{\subsection}{0pt}{6pt}{2pt}
\titlespacing{\subsubsection}{0pt}{4pt}{2pt}

% ── Caption style ────────────────────────────────────────────
\captionsetup{
  font={small,it},
  labelfont={bf,color=inkmuted},
  labelsep=period,
  skip=3pt,
}

% ── List style ───────────────────────────────────────────────
\setlist[itemize]{
  leftmargin=*,
  topsep=2pt,
  itemsep=1pt,
  parsep=0pt,
  label={\textcolor{etred}{\textbullet}},
}

% ── Inline code style ────────────────────────────────────────
\newcommand{\code}[1]{{\ttfamily\small\color{steel}#1}}

% ── Formula box ─────────────────────────────────────────────
\newmdenv[
  linecolor=etgold,
  linewidth=1pt,
  backgroundcolor=white,
  innerleftmargin=6pt,
  innerrightmargin=6pt,
  innertopmargin=5pt,
  innerbottommargin=5pt,
  skipabove=4pt,
  skipbelow=4pt,
]{formulabox}

% ── Red rule box (highlight) ─────────────────────────────────
\newmdenv[
  linecolor=etred,
  linewidth=1.5pt,
  backgroundcolor=reddim,
  innerleftmargin=6pt,
  innerrightmargin=6pt,
  innertopmargin=4pt,
  innerbottommargin=4pt,
  skipabove=4pt,
  skipbelow=4pt,
]{redbox}

% ── Header / Footer ─────────────────────────────────────────
\pagestyle{fancy}
\fancyhf{}
\renewcommand{\headrulewidth}{0pt}
\fancyhead[L]{\footnotesize\color{inkmuted}\textit{EthosNews — ET GenAI Hackathon 2026}}
\fancyhead[R]{\footnotesize\color{inkmuted}\textit{Phase 2 Prototype}}
\fancyfoot[C]{\footnotesize\color{inkmuted}\thepage}

% ── Paragraph spacing ────────────────────────────────────────
\setlength{\parskip}{3pt}
\setlength{\parindent}{0pt}

% ============================================================
\begin{document}

% ── MASTHEAD (full width) ────────────────────────────────────
\twocolumn[{%
  \begin{center}
    {\color{etred}\rule{\textwidth}{3pt}}\\[4pt]
    {\fontfamily{lmr}\selectfont\Huge\bfseries\color{ink} ETHOSNEWS}\\[2pt]
    {\large\color{inkmuted}\textit{From Algorithmic Feeds to Epistemic Agency}}\\[3pt]
    {\small\color{inkmuted}%
      \textbf{Economic Times GenAI Hackathon 2026} \quad\textcolor{rule}{|}\quad
      Phase 2 Prototype Submission \quad\textcolor{rule}{|}\quad
      March 2026}\\[4pt]
    {\color{etred}\rule{\textwidth}{1pt}}\\[2pt]
    {\color{rule}\rule{\textwidth}{0.5pt}}\\[8pt]
  \end{center}
}]

% ============================================================
\section{Product Overview}

In 2026, business news is still delivered as it was two decades ago — static text, one feed for everyone, no way to know what to trust. EthosNews is a \textbf{participatory, intent-driven news ecosystem} built to fix that. It places the user's own values, interests, and epistemic preferences at the centre of every decision the system makes — from which articles surface, to how verified each one is, to what other perspectives exist on the same story.

\begin{redbox}
\small\textbf{Core idea:} A user's intent is not a filter applied on top of an algorithm. It \textit{is} the algorithm.
\end{redbox}

\subsection{The Three Problems}

\textbf{Cognitive Capture.} Engagement algorithms repeatedly surface 3–5 dominant viewpoints, creating a narrative monoculture. A user reads widely but sees narrowly.

\textbf{Information Overload.} Without semantic filtering tied to a reader's actual goals, volume becomes noise. Most articles are irrelevant to most readers, yet arrive anyway.

\textbf{Truth Decay.} AI-generated content has degraded trust signals. Verification lag means false claims spread faster than corrections. Readers have lost calibration for what to believe.

\subsection{ET Hackathon Tracks Addressed}

\textbf{My ET — The Personalised Newsroom.} A mutual fund investor, a startup founder, and a student each receive a structurally different news experience — driven not by topic tags alone but by epistemic mode, verification threshold, diversity preference, and reading depth.

\textbf{News Navigator — Interactive Intelligence Briefings.} Instead of reading eight separate articles on one story, a user gets a single AI-powered layer: claim-by-claim verification with evidence, and a live map of every alternative perspective on that story.

\subsection{Who It Is For}

\begin{center}
\small
\begin{tabularx}{\columnwidth}{@{}lX@{}}
\toprule
\textbf{User} & \textbf{How EthosNews Serves Them} \\
\midrule
General Consumer & PNC-filtered feed + AI slop score on every article \\
Journalist & One-click fact-check + narrative cluster map \\
Fact-Checker & 4-stage pipeline: extract → gather → rerank → classify \\
Student & Readability depth + data density controls in PNC \\
\bottomrule
\end{tabularx}
\end{center}

\subsection{The Daily Experience}

A user opens EthosNews and immediately sees their \textbf{Personalised Feed} — articles ranked by how well they match a constitution the user defined once, in plain English. On any article, three AI actions are one click away: \textbf{Verify Claims} runs a real-time fact-check; \textbf{Narrative Clusters} surfaces every alternative perspective on the story; \textbf{Listen} generates a narrated audio version in the user's chosen style. Every article also carries a \textbf{Slop Meter} — a locally-computed score indicating the likelihood that the content was AI-generated, with no external API required.

% ============================================================
\section{System Architecture}

EthosNews has four layers: a React frontend, a FastAPI backend, an AI processing layer, and a three-database data layer. All communication uses JWT-authenticated REST.

\begin{figure}[H]
  \centering
  \includegraphics[width=\columnwidth,height=7cm,keepaspectratio]{fig_arch.png}
  \caption{End-to-end system architecture.}
\end{figure}

\subsection{Technology Stack}

\begin{center}
\small
\begin{tabularx}{\columnwidth}{@{}lX@{}}
\toprule
\textbf{Layer} & \textbf{Technology} \\
\midrule
Frontend      & React 19, Vite 8, framer-motion \\
Backend       & Python 3.11, FastAPI, uv \\
Primary DB    & PostgreSQL + asyncpg \\
Cache         & Redis (TTL: 2–60 min per type) \\
Vector DB     & Qdrant (cosine, 384-dim) \\
Embeddings    & SentenceTransformer all-MiniLM-L6-v2 \\
Reranking     & CrossEncoder ms-marco-MiniLM-L-6-v2 \\
NLP           & spaCy en\_core\_web\_sm \\
LLM (fast)    & llama-3.1-8b-instant / gpt-4o-mini \\
LLM (strong)  & llama-3.3-70b-versatile / gpt-4o \\
Auth          & JWT HS256, bcrypt \\
Data source   & GDELT (real-time global news stream) \\
\bottomrule
\end{tabularx}
\end{center}

\subsection{Ingestion Pipeline}

Every 15 minutes, a background process polls GDELT for new articles. What arrives is a stream of event metadata and source URLs — the actual article text has to be fetched, cleaned, and understood before it is useful.

\begin{figure}[H]
  \centering
  \includegraphics[width=\columnwidth,height=8cm,keepaspectratio]{fig_ingest.png}
  \caption{GDELT ingestion pipeline — producer/consumer with 3 parallel workers.}
\end{figure}

\textbf{Clever detail — deterministic IDs.} Each article's database ID is computed as \code{uuid5(NAMESPACE\_URL, url)}. The same URL always produces the same UUID. This makes deduplication a simple point-existence check in Qdrant — no separate seen-set, no extra lookup table.

Articles shorter than 150 words are dropped. Survivors are scraped, entity-extracted via spaCy, embedded into a 384-dimensional vector, scored for AI generation, and stored simultaneously in Qdrant (for semantic search) and PostgreSQL (for structured queries).

% ============================================================
\section{Personal News Constitution}

The Personal News Constitution (PNC) is the identity layer of EthosNews. It is not a list of topic preferences — it is a structured representation of \textit{how} a user relates to information, not just \textit{what} they want to read about.

\subsection{The Schema}

A user types a description in plain English. A strong-tier LLM (temp=0.2) reads it and extracts a four-field structured document:

\begin{center}
\small
\begin{tabularx}{\columnwidth}{@{}lX@{}}
\toprule
\textbf{Field} & \textbf{What It Controls} \\
\midrule
\texttt{epistemic\_framework} & Primary mode (empiricist, rationalist, narrative, pragmatist, humanist) + verification threshold 0–1 \\
\texttt{narrative\_preferences} & Diversity weight 0–1 + bias tolerance (low/medium/high) \\
\texttt{topical\_constraints} & Priority domains (max 10) + excluded topics \\
\texttt{complexity\_preference} & Readability depth (beginner/intermediate/expert) + data density \\
\bottomrule
\end{tabularx}
\end{center}

\texttt{epistemic\_framework.primary\_mode} is the most novel field. It does not describe \textit{what} topics the user wants — it describes \textit{how} they want to relate to information. An \textit{empiricist} wants high-evidence data-backed stories. A \textit{narrative} reader wants rich contextual storytelling. This mode propagates into LLM prompts throughout the system. If the LLM call fails, hardcoded defaults are used so the user always gets a feed.

\subsection{How PNC Drives the Feed}

\begin{figure}[H]
  \centering
  \includegraphics[width=\columnwidth,height=7cm,keepaspectratio]{fig_pnc.png}
  \caption{From natural-language constitution to personalised feed.}
\end{figure}

The user's \texttt{priority\_domains} are joined into a string and encoded into a 384-dim query vector. Qdrant returns the top 200 semantically similar articles. Each candidate is then scored with a time-decay function:

\begin{formulabox}
\small
\centering
$\text{score} = \text{similarity} \times 0.85^{\,\text{days\_old}}$
\end{formulabox}

A 3-day-old article retains 61\% of its semantic score; a 7-day-old article retains 32\%. This keeps the feed fresh without discarding relevant older reporting. Articles matching \texttt{excluded\_topics} are removed by substring filter. The top 30 survivors are returned.

% ============================================================
\section{AI Pipelines}

\subsection{Agentic Fact-Checking}

When a user clicks \textit{Verify Claims}, a 4-stage pipeline runs on the article in real time.

\begin{figure}[H]
  \centering
  \includegraphics[width=\columnwidth,keepaspectratio]{img_factcheck.png}
  \caption{Fact-check pipeline: extract → evidence → rerank → classify.}
\end{figure}

\subsubsection{Stage 1 — Claim Extraction}
A fast-tier LLM (temp=0.0) breaks the article into at most 8 atomic factual claims — statements that can, in principle, be verified against an external source. If the LLM fails, the system falls back to simple sentence splitting.

\subsubsection{Stage 2 — Evidence Gathering (Parallel)}
For every claim, two searches run simultaneously: a semantic search against the Qdrant article corpus (top 15 hits), and a live web search via Tavily (top 5 results, only if an API key is configured). Each result is trimmed to 600 characters of context.

\subsubsection{Stage 3 — Batched CrossEncoder Reranking}
All \texttt{[claim, passage]} pairs across \textit{all} claims are flattened into a single list and scored in one \code{ce.predict()} call. This batching is the most computationally intensive step — doing it in one forward pass rather than N×M individual calls is a significant latency saving. The top 4 passages per claim are kept. Raw scores are normalised via sigmoid with a softening divisor of 3.0, preventing moderate-confidence evidence from collapsing to 0 or 1.

\subsubsection{Stage 4 — LLM Classification}
Each claim is judged against its top-4 evidence passages by a fast-tier LLM. Concurrent calls are throttled by \code{asyncio.Semaphore(3)} to respect API rate limits. Each claim receives one of three verdicts: \textbf{Supported}, \textbf{Contradicted}, or \textbf{Not Mentioned}, along with a confidence score and supporting URLs. The fraction of \textit{Not Mentioned} claims is surfaced to the user as the \textbf{unverifiable ratio} — a proxy for how well-corroborated the story is.

\subsection{Narrative Divergence Clustering}

When a user clicks \textit{Narrative Clusters}, the system maps every alternative perspective on the article's story — not just similar articles, but the full landscape of how the same event is being framed differently.

\begin{figure}[H]
  \centering
  \includegraphics[width=\columnwidth,keepaspectratio]{img_clustering.png}
  \caption{Narrative divergence: HDBSCAN groups articles by semantic distance into viewpoint clusters.}
\end{figure}

Qdrant retrieves the 50 most semantically similar articles. HDBSCAN clusters them by their 384-dim vectors (\code{min\_cluster\_size=3}, Euclidean distance, EOM selection). Each cluster is a distinct \textbf{narrative pillar} — a coherent viewpoint on the story.

Every pillar receives a \textbf{divergence score}:

\begin{formulabox}
\small
\centering
$\text{divergence} = 0.8 \times \text{cosine\_dist} + 0.2 \times \text{tone\_std}$
\end{formulabox}

\texttt{cosine\_dist} measures how far the cluster's articles sit from its own centroid — high distance means internally diverse viewpoints. \texttt{tone\_std} is the standard deviation of GDELT's pre-computed sentiment scores across the cluster — high variance means the cluster mixes positive and negative framings of the same event, a signal of genuine narrative tension not visible in semantic space alone.

Three representative article URLs per cluster are selected by blending centroid proximity (weight 0.6) and entity coverage breadth (weight 0.4). A strong-tier LLM (temp=0.3) writes a 2–3 sentence plain-English summary for each pillar. HDBSCAN is used rather than k-means because it does not require specifying the number of clusters in advance and handles noise points gracefully — important when story clusters are irregular in size.

\subsection{Voice Narration}

Any article can be listened to in three broadcast styles, generated end-to-end without pre-recorded assets.

\begin{center}
\small
\begin{tabularx}{\columnwidth}{@{}llX@{}}
\toprule
\textbf{Mode} & \textbf{Voice} & \textbf{Style} \\
\midrule
Anchor  & Fritz-PlayAI  & Authoritative broadcast — lead with the single most important fact \\
Podcast & Nova-PlayAI   & Conversational, warm, rhetorical questions, like explaining to a friend \\
Drama   & Thunder-PlayAI & Cinematic tension, short punchy sentences, gut-punch close \\
\bottomrule
\end{tabularx}
\end{center}

A fast-tier LLM (temp=0.75) generates a \textasciitilde100-word spoken script tailored to the mode's system prompt. The Groq TTS model (\code{playai-tts}) synthesises it to MP3. Because the Groq SDK is synchronous, the call runs inside \code{asyncio.to\_thread()} to avoid blocking the FastAPI event loop. The response returns the script text alongside a base64-encoded MP3 — the frontend plays it directly without a file download.

\subsection{AI Slop Detection}

Every ingested article is scored for AI-generation likelihood using five statistical signals computed entirely locally — no LLM call, no external API.

\begin{center}
\small
\begin{tabularx}{\columnwidth}{@{}Xlr@{}}
\toprule
\textbf{Signal} & \textbf{What It Measures} & \textbf{Weight} \\
\midrule
Burstiness        & Sentence-length variation (AI text is uniform) & 0.30 \\
N-gram Repetition & Repeated 4-gram frequency                      & 0.25 \\
Entity Density    & Named entities per 100 words                   & 0.20 \\
Passive Voice     & Fraction of passive sentences                  & 0.15 \\
Uniformity        & \% of sentences in 12–25 word range            & 0.10 \\
\bottomrule
\end{tabularx}
\end{center}

\begin{formulabox}
\small
\centering
$\text{score} = 0.30\,B + 0.25\,R + 0.20\,E + 0.15\,P + 0.10\,U$

\smallskip
$< 0.35 \rightarrow \textbf{human} \quad
 0.35\text{–}0.65 \rightarrow \textbf{uncertain} \quad
 > 0.65 \rightarrow \textbf{ai\_generated}$
\end{formulabox}

Only texts scoring \textit{uncertain} trigger a Layer 2 LLM review, keeping costs and latency low for clear cases. The merged score weights L1 at 0.4 and L2 at 0.6 — when statistical signals are inconclusive, semantic judgement is more informative.

% ============================================================
\section{Leaderboard \& Gamification}

The leaderboard redefines what good engagement looks like. A user earns more not by spending more time on the platform, but by consuming content that is constitutionally aligned, perspectivally diverse, and verified.

\subsection{Per-Read Engagement Score}

Every article read generates an \code{engagement\_event} with three components:

\begin{formulabox}
\small
\centering
$\text{score} = 0.40\,\text{PNC\_align} + 0.30\,\text{diversity} + 0.30\,\text{time}$

\smallskip
$\text{time} = \dfrac{\min(\text{secs}, 600)}{600} \times (1 - \text{slop\_score})$
\end{formulabox}

\textbf{PNC alignment} is the cosine similarity between the article's embedding and the user's constitution vector, normalised to [0,1]. \textbf{Diversity} is 1.0 if the article spans two or more narrative clusters, 0.5 otherwise — rewarding users who read across perspectives. \textbf{Time} is capped at 10 minutes and \textit{down-weighted by the article's slop score} — spending time on AI-generated content accrues less reward than the same time on verified human journalism.

\subsection{Ranking Formula}

\begin{formulabox}
\small
\centering
$\text{rank} = (\text{reads} \times 1.0 + \text{participations} \times 2.5)$
$\times\;(1 + \log_{10}(\text{streak}+1) \times 0.5)$
\end{formulabox}

Active participation (triggering fact-checks, exploring clusters) carries 2.5× the weight of passive reading. The \(\log_{10}\) streak multiplier gives diminishing returns — a long-time user's advantage over a new engaged user is bounded, preventing compounding lock-in.

All events are written to an immutable \code{engagement\_events} table in PostgreSQL (full audit trail with component breakdown). The running leaderboard total is maintained in a Redis sorted set — O(log N) inserts, O(log N + K) top-K reads, with atomic increments.

% ============================================================
\section{Implementation}

\subsection{Backend Services}

\begin{center}
\small
\begin{tabularx}{\columnwidth}{@{}lX@{}}
\toprule
\textbf{File} & \textbf{Responsibility} \\
\midrule
\texttt{core/clients.py}       & Thread-safe singleton registry for all heavy ML models — SentenceTransformer, CrossEncoder, spaCy, Qdrant. Double-checked locking prevents redundant initialisation under concurrent requests. \\
\texttt{core/llm.py}           & Unified async \texttt{chat()} interface over Groq/OpenAI/Gemini. Dual-tier (fast/strong). Gemini's sync SDK runs in \texttt{run\_in\_executor}. \\
\texttt{core/cache.py}         & Redis helpers with per-type TTLs: feed 2–5 min, articles 1 hr, fact-checks 1 hr, clusters 30 min. \\
\texttt{services/ingestion/}   & GDELT producer-consumer. \texttt{uuid5} dedup. Scrape → embed → slop → store. \\
\texttt{services/fact\_checker/} & 4-stage pipeline. Batched CrossEncoder in single forward pass. \\
\texttt{services/clustering/}  & HDBSCAN on Qdrant vectors. Divergence = cosine + tone\_std blend. \\
\texttt{services/slop\_detector/} & L1 statistical (5 signals). L2 LLM fallback only for \textit{uncertain} band. \\
\texttt{services/pnc\_service.py} & LLM extraction with JSON schema injection. Hardcoded fallback defaults. \\
\texttt{services/leaderboard/} & Per-read scoring. Weekly aggregate with log streak multiplier. Redis + Postgres dual write. \\
\texttt{api/voice.py}          & LLM script (3 modes) + Groq TTS in \texttt{asyncio.to\_thread}. \\
\texttt{api/feed.py}           & \texttt{\_excerpt\_from\_text()} strips URLs + whitespace, truncates to 320 chars. Regex pre-compiled at module load. \texttt{\_serialize\_qdrant\_hit()} normalises raw Qdrant payload for feed responses. \\
\bottomrule
\end{tabularx}
\end{center}

\subsection{API Surface}

\begin{center}
\small
\begin{tabularx}{\columnwidth}{@{}lX@{}}
\toprule
\textbf{Router} & \textbf{Endpoints} \\
\midrule
Auth        & \texttt{POST /auth/signup} · \texttt{login} · \texttt{GET /auth/me} \\
Feed        & \texttt{GET /feed} · \texttt{/article/\{id\}} · \texttt{/personalized\_feed/\{uid\}} \\
Fact-Check  & \texttt{POST /article/\{id\}/fact-check} · \texttt{/fact-check} \\
Clusters    & \texttt{GET /article/\{id\}/clusters} \\
PNC         & \texttt{POST /pnc/generate} · \texttt{/pnc/\{uid\}} · \texttt{PATCH /pnc/\{uid\}} \\
Leaderboard & \texttt{POST /leaderboard/event} · \texttt{GET /leaderboard/top} \\
Voice       & \texttt{GET /article/\{id\}/voice?mode=anchor|podcast|drama} \\
\bottomrule
\end{tabularx}
\end{center}

\subsection{Database Schema}

\textbf{users} — id, username, email, password\_hash, onboarding\_completed, streak\_count, active\_participations

\textbf{articles} — id (\textit{uuid5 of url}), title, content, url, source, image\_url, published\_at, category, avg\_tone, num\_mentions, country\_code, ai\_slop\_score, ai\_slop\_label

\textbf{user\_constitutions} — user\_id (PK/FK), constitution (JSONB), updated\_at

\textbf{engagement\_events} — id, user\_id, article\_id, pnc\_alignment, diversity\_score, time\_score, total\_score, created\_at

% ============================================================
\section{Design Decisions}

\begin{description}[style=unboxed, leftmargin=0pt, itemsep=3pt]

\item[\textcolor{etred}{Qdrant over FAISS.}]
Qdrant supports payload filtering \textit{alongside} vector search. Category-scoped cluster retrieval works without pulling all vectors into memory first.

\item[\textcolor{etred}{HDBSCAN over k-means.}]
Story clusters are irregular in size and shape. HDBSCAN requires no \textit{k} in advance and handles noise points gracefully — both critical for open-ended news clustering.

\item[\textcolor{etred}{uuid5 for article IDs.}]
Deterministic IDs from URLs make deduplication idempotent. Re-ingesting the same article is a no-op with a single Qdrant point-existence check.

\item[\textcolor{etred}{GDELT \texttt{avg\_tone} in divergence scoring.}]
GDELT provides pre-computed sentiment per event. Using it as a divergence signal adds a non-embedding dimension that captures framing differences not visible in pure semantic space.

\item[\textcolor{etred}{Redis sorted set for leaderboard.}]
O(log N) writes and reads, atomic increments, natural top-K queries — no aggregation on Postgres for real-time leaderboard reads.

\item[\textcolor{etred}{Slop L2 only on uncertain band.}]
Running an LLM on every article would be expensive. Statistical signals handle the clear cases cheaply; the LLM is reserved for the ambiguous middle, keeping p95 latency low.

\item[\textcolor{etred}{TTS in thread pool.}]
The Groq TTS SDK is synchronous. \code{asyncio.to\_thread()} offloads the blocking call to a worker thread so the event loop stays responsive during audio generation.

\end{description}

% ============================================================
\section{Future Directions}

\begin{itemize}
  \item \textbf{Rephrase API} — adapt article content per PNC \texttt{readability\_depth} and \texttt{epistemic\_framework}: the same story rewritten for a beginner versus an expert, or for a narrative reader versus an empiricist.
  \item \textbf{PNC feedback loop} — implicit signals from reading behaviour (time spent, clusters explored, fact-checks triggered) fed back to refine the constitution over time.
  \item \textbf{Vernacular adaptation} — culturally-adapted translation into Hindi, Tamil, Telugu, Bengali using the existing LLM provider abstraction — not literal translation but context-aware localisation.
  \item \textbf{Browser extension} — PNC-based slop detection and claim verification as an overlay on any news website, without opening the app.
  \item \textbf{LLM-as-Judge moderation} — evaluating user comment quality against constitutional alignment, distinguishing substantive contribution from noise.
  \item \textbf{Mobile application} — native iOS/Android with push alerts for high-divergence story events.
\end{itemize}

% ── Footer rule ──────────────────────────────────────────────
\vfill
\noindent{\color{etred}\rule{\columnwidth}{1.5pt}}\\[2pt]
\noindent{\footnotesize\color{inkmuted}
  \textit{EthosNews Phase 2 Prototype} \quad\textcolor{rule}{|}\quad
  \textit{ET GenAI Hackathon 2026} \quad\textcolor{rule}{|}\quad
  \textit{Full source: FastAPI + React 19 + Qdrant + PostgreSQL + Redis}
}

\end{document}
